\section{What the ladder dataset shows}
\label{sec:offline}

\subsection{Regret is rare and concentrated}

Table~\ref{tab:uniform} gives the price of a uniform budget. At 200 visits the engine gives up 0.99\% of winrate per move on average, and doubling the budget removes roughly a third of what remains. The average hides the shape. At 200 visits 71\% of positions carry no regret at all: the 200-visit move is the best move we scored. The worst tenth of positions hold 84\% of the total. Regret also depends on when in the game the position occurs. At 200 visits it averages 0.19\% per move over the first 40 moves, 0.74\% from move 40 to 100, and 1.4 to 1.6\% after move 100. The opening is nearly free of search error at these budgets, which will matter for how we interpret memory.

A rule that knew each position's regret and stopped at the first rung with none would match uniform search at twelve times the visits at a 200-visit budget (last column of Table~\ref{tab:frontier}). That is the room available to any stopping rule on this ladder.

\begin{table}[t]
\centering\small
\begin{tabular}{rrrrr}
\toprule
Visits & Mean regret & Positions with & Regret held by & Move differs from \\
 & (\% winrate) & no regret (\%) & worst 10\% (\%) & 3{,}200-visit move (\%) \\
\midrule
50 & 1.47 & 66 & 78 & 33 \\
100 & 1.19 & 68 & 81 & 30 \\
200 & 0.99 & 71 & 84 & 27 \\
400 & 0.70 & 74 & 86 & 22 \\
800 & 0.48 & 78 & 90 & 17 \\
1{,}600 & 0.30 & 82 & 94 & 11 \\
3{,}200 & 0.17 & 86 & 98 & 0 \\
\bottomrule
\end{tabular}

\caption{Uniform search on the ladder dataset (6{,}000 positions). Regret is measured against the best move scored by an independent forced search.}
\label{tab:uniform}
\end{table}

\subsection{No single signal finds it}

Table~\ref{tab:auc} asks how well individual search statistics rank positions by whether more than 0.5\% of winrate remains on the table. The measures of how contested the root is (top move's visit share, the margin between the best move's lower confidence bound and its best rival) reach an AUC of about 0.74. KataGo's own prediction of its short-term winrate error, a signal it uses inside the tree to weight playouts, reaches 0.58 to 0.60 as a predictor of whether the root decision will improve. The classical time-management cue, that the most-visited move is not the highest-valued one \citep{huang2010time}, is weaker still at 0.59 to 0.61. Contested positions are common, and most of them are contested between moves of nearly equal value, where being wrong costs nothing.

\begin{table}[t]
\centering\small
\begin{tabular}{lrrr}
\toprule
Signal & 50 & 200 & 800 \\
\midrule
Top move's visit share & 0.75 & 0.74 & 0.74 \\
LCB margin & 0.74 & 0.73 & 0.73 \\
v1 entropy gate & 0.73 & 0.73 & 0.74 \\
KL(visits, prior) & 0.65 & 0.63 & 0.61 \\
Most-visited is not best-valued & 0.60 & 0.59 & 0.61 \\
KataGo's predicted winrate error & 0.58 & 0.58 & 0.60 \\
Searched minus raw winrate & 0.57 & 0.56 & 0.56 \\
\bottomrule
\end{tabular}

\caption{AUC for ranking positions by remaining regret above 0.5\%, from the search at 50, 200, and 800 visits.}
\label{tab:auc}
\end{table}

\subsection{A learned stop on a ladder beats uniform search by 1.7 to 1.9 times}

Table~\ref{tab:frontier} is the main offline result. With the 50, 200, 800, 3{,}200 ladder and a mean cost of 200 visits, the learned stopper gives up 0.75\% of winrate per move. Uniform search needs 358 visits to do as well, a multiplier of 1.80. The multiplier stays between 1.69 and 1.88 for mean budgets from 100 to 800 visits. If every restarted search is charged in full, so that reaching the 800-visit rung costs $50+200+800$, the multiplier is 1.29 to 1.51.

\begin{table}[t]
\centering\small
\begin{tabular}{rrrrrrr}
\toprule
Mean & Uniform & Stopper & Equivalent & \multicolumn{2}{c}{Multiplier} & Oracle \\
visits & regret (\%) & regret (\%) & uniform visits & continuation & restart & multiplier \\
\midrule
100 & 1.19 & 1.03 & 173 & 1.76 & 1.51 & 5.5 \\
150 & 1.07 & 0.87 & 265 & 1.79 & 1.47 & 9.6 \\
200 & 0.99 & 0.75 & 358 & 1.80 & 1.44 & 12.1 \\
300 & 0.82 & 0.60 & 544 & 1.88 & 1.47 & 10.9 \\
400 & 0.70 & 0.53 & 692 & 1.75 & 1.36 & 8.1 \\
600 & 0.57 & 0.41 & 1072 & 1.80 & 1.38 & 5.7 \\
800 & 0.48 & 0.35 & 1332 & 1.69 & 1.29 & 5.7 \\
\bottomrule
\end{tabular}

\caption{Learned stopping on the 50/200/800/3{,}200 ladder, cross-fitted by game. \emph{Equivalent uniform visits} is the fixed budget with the same mean regret. The oracle stops with knowledge of the true regret.}
\label{tab:frontier}
\end{table}

\subsection{The ladder and the rate rule do most of the work}

Table~\ref{tab:rules} runs hand-written stopping signals through the same ladder. Two things stand out. First, a flat threshold is harmful for most signals: v1's entropy gate with a flat threshold reaches a multiplier of 0.63 at 200 visits, which is worse than not gating at all, and random stopping is far worse. \citet{lan2021learning} made the same observation about random stopping. Second, dividing the score by the visits needed to reach the next rung turns the same signals into gains. The LCB margin alone, with no learning, reaches 1.45 at 200 visits.

\begin{table}[t]
\centering\small
\begin{tabular}{lrrrr}
\toprule
 & \multicolumn{2}{c}{Flat threshold} & \multicolumn{2}{c}{Rate rule} \\
Stopping signal & 200 visits & 400 visits & 200 visits & 400 visits \\
\midrule
Random stopping & 0.44 & 0.29 & 0.71 & 0.69 \\
Most-visited is not best-valued & 0.55 & 0.60 & 0.58 & 0.77 \\
v1 entropy gate & 0.63 & 0.76 & 1.16 & 1.28 \\
LCB margin & 1.13 & 0.86 & 1.45 & 1.00 \\
\midrule
Learned stopper & 1.41 & 1.22 & 1.82 & 1.77 \\
\bottomrule
\end{tabular}

\caption{Compute multipliers (continuation cost) for stopping signals on the 50/200/800/3{,}200 ladder. Values below 1 are worse than uniform search.}
\label{tab:rules}
\end{table}
